% ── Übung: Sitz ────────────────────────────────────────────────────────
\begin{uebung}{Sitz}{welpen}{Grundkommando}{2}
    \uebungsbild{images/heimuebungen/welpen/sitz.png}

    \uebungsbeschreibung{%
        Der Hund lernt sich auf ein Sicht- und/oder Wortsignal zu setzen.\\[6pt]
        \textbf{Schritte:}
        \begin{enumerate}[leftmargin=*, itemsep=2pt]
      \item Ein Leckerli wird in der Hand an die Nase des Hundes geführt. Der Zeigefinger ist dabei nach oben gestreckt.
      \item Sobald der Hund das Leckerli bemerkt hat und sich daran orientiert die Hand nach oben, leicht über den Kopf des Hundes führen.
      \item Wenn der Hund sich setzt das Verhalten \textit{markieren} (Clicker oder Wort) und belohnen.
      \item Das Signalwort erst einführen, wenn das Verhalten zuverlässig gezeigt wird.
            \item Wichtig: Auflösungskommando bevor der Hund selber den Befehl auflöst.
        \end{enumerate}
    }

    \uebungsvariationen{%
        \begin{itemize}[leftmargin=*, itemsep=4pt]
      \item \textbf{Distanz:} Sitz mit mehr Abstand üben, so dass das Kommando Sitz auch aus der Entfernung klappt.
      \item \textbf{Ablenkung steigern:} Wenn es zu Hause gut funktioniert in den Garten, nach draußen oder auch an belebtere Orte verlagern.
      \item \textbf{Dauer steugern:} Der Hund soll lernen die Sitzposition länger zu halten. Das Auflösungswort herauszögern.
        \end{itemize}
    }

    \uebungstrainer{%
        \begin{itemize}[leftmargin=*, itemsep=4pt]
      \item Hand nicht zu hoch halten, so dass der KHund springt.
      \item DAuf das Timing bei der Best�ätigung achten.
      \item Auf Auflösewort achten!
        \end{itemize}
    }

\end{uebung}
